% =====================================================================
%  ANEXO A — Cuadros comparativos del Marco Teórico
%  (trasladados desde el Capítulo II para cumplir el límite de 10 páginas)
% =====================================================================
\chapter{CUADROS COMPARATIVOS DEL MARCO TEÓRICO}
\label{anx:comparativas}

Este anexo reúne los cuadros comparativos que sustentan las decisiones
tecnológicas y metodológicas descritas en el Capítulo~\ref{cap:marco}.

{\small
\begin{longtable}{|L{2.6cm}|L{6.65cm}|L{5.4cm}|}
\hline
\rowcolor{encabezado} \textbf{Metodología} & \textbf{Descripción} & \textbf{Características} \\
\hline
\endfirsthead
\hline
\rowcolor{encabezado} \textbf{Metodología} & \textbf{Descripción} & \textbf{Características} \\
\hline
\endhead

\endfoot
\textbf{Waterfall} & La metodología waterfall o cascada, es un proceso
secuencial, el proyecto se divide en fases, las cuales deben completarse
completamente antes de pasar a la siguiente. El avance del progreso va de
arriba hacia abajo. &
\begin{tablista}
  \item Enfoque lineal y estructurado.
  \item Requiere una planificación previa.
  \item No permite adaptaciones una vez iniciada la ejecución.
  \item Fácil de entender.
\end{tablista} \\
\hline
\textbf{Espiral} & Esta metodología integra la gestión de riesgos a lo largo
del proceso de desarrollo. El desarrollo ocurre en ciclos o vueltas y cada uno
de estos cuenta con 4 fases: planificación, análisis de riesgos, ingeniería y
evaluación por parte del cliente. &
\begin{tablista}
  \item Enfoque iterativo basado en ciclos sucesivos.
  \item Permite incorporar ajustes o mejoras en cada iteración.
  \item Contempla análisis de riesgos periódico.
\end{tablista} \\
\hline
\textbf{Incremental} & Esta metodología combina elementos de la metodología
cascada, permite desarrollar el producto de forma progresiva, cada etapa
incremental constituye una nueva funcionalidad, lo que permite ver resultados
de forma más rápida en comparación del modelo cascada. &
\begin{tablista}
  \item Entrega parcial del sistema.
  \item Facilita la detección de errores temprana.
  \item Recopila comentarios de clientes después de cada entrega.
  \item Flexible a los cambios.
  \item Combinación de enfoques lineales e iterativos.
\end{tablista} \\
\hline
\textbf{Prototipado} & Esta metodología desarrolla un modelo preliminar y
tangible de software, esto permite que el usuario pruebe y de
retroalimentación, para así poder ajustar los requisitos y mejorar el diseño
del sistema final. &
\begin{tablista}
  \item Facilita la comunicación con el usuario.
  \item Reduce la incertidumbre.
  \item Es iterativa.
  \item Reduce riesgos y costos.
  \item Es flexible y versátil.
\end{tablista} \\
\hline
\textbf{Diseño rápido de aplicaciones RAD} & Esta metodología permite
desarrollar software de alta calidad en un corto periodo de tiempo. Los costes
son mucho más altos y el desarrollo es flexible, aunque requiere mayor
intervención de los usuarios. &
\begin{tablista}
  \item Desarrollo basado en prototipos funcionales.
  \item Iteraciones cortas y desarrollo incremental.
  \item Alta participación del usuario.
\end{tablista} \\ \hline
\caption[Comparación de Metodologías tradicionales de desarrollo de software]{Comparación de Metodologías tradicionales de desarrollo de software. \textit{Fuente:} Elaboración propia.}\label{tab:metodologias-tradicionales}\\
\end{longtable}
}

{\small
\begin{longtable}{|L{2.4cm}|L{7.05cm}|L{5.2cm}|}
\hline
\rowcolor{encabezado} \textbf{Metodología} & \textbf{Descripción} & \textbf{Características} \\
\hline
\endfirsthead
\hline
\rowcolor{encabezado} \textbf{Metodología} & \textbf{Descripción} & \textbf{Características} \\
\hline
\endhead

\endfoot
\textbf{Scrum} & Es una metodología ágil, que divide los requerimientos y
tareas de forma similar a Kanban. Está compuesto por iteraciones fijas y cortas
denominadas Sprints, así como una estructura basada en roles, eventos y
artefactos definidos. Su objetivo es entregar incrementos funcionales del
producto en ciclos de 2 a 4 semanas, fomentando la colaboración y la
inspección continua (Schwaber \& Sutherland, 2020). &
\begin{tablista}
  \item Roles definidos (Product Owner, Scrum Master, Equipo).
  \item Iteraciones cortas y fijas.
  \item Entregas incrementales al final de cada sprint.
  \item Alta colaboración con el cliente.
  \item Uso de backlog y tableros para la gestión visual del trabajo.
  \item[] (Schwaber \& Sutherland, 2020)
\end{tablista} \\
\hline
\textbf{Extreme Programming (XP)} & Es una metodología basada en relaciones
interpersonales, la colaboración y la retroalimentación continua. Su objetivo
principal es crear un entorno de trabajo basado en la comunicación constante,
simplicidad del diseño y participación activa del cliente. Incorpora prácticas
de ingeniería como pair programming, test-driven development, refactorización
constante e integración continua (Beck \& Andres, 2005). &
\begin{tablista}
  \item Trabajo colaborativo.
  \item Diseño sencillo y adaptables.
  \item Desarrollo guiado por pruebas.
  \item Refactorización continua.
  \item Integración continua.
  \item Entregas frecuentes y pequeñas iteraciones.
  \item Planificación basada en historias de usuario.
  \item Programación en parejas.
\end{tablista} \\
\hline
\textbf{Mobile-D} & Es una metodología ágil para el desarrollo de aplicaciones
móviles para equipos pequeños, está centrada en ciclos de desarrollo rápidos y
entrega continua. Combina prácticas ágiles con procesos orientados al hardware
y ciclos muy cortos de entrega (Abrahamsson, y otros, 2004). &
\begin{tablista}
  \item Diseñada para desarrollo móvil.
  \item Iteraciones rápidas.
  \item Adaptación a restricciones de hardware.
  \item Enfocado en grupos pequeños.
\end{tablista} \\
\hline
\textbf{Kanban} & Es una metodología de trabajo inventada por la empresa de
automóviles Toyota. Consiste en dividir las tareas en porciones mínimas y
organizarlas en un tablero de trabajo dividido en tareas pendientes, en curso y
finalizadas. De esta forma se crea un flujo de trabajo muy visual basado en
tareas prioritarias e incrementando el valor del producto (Santander, 2020). &
\begin{tablista}
  \item Visualización del trabajo.
  \item Límite de trabajo en curso.
  \item Mejora continua.
  \item No requiere iteraciones fijas.
  \item Fácil implementación en equipos individuales.
  \item Gestión visual del trabajo.
\end{tablista} \\ \hline
\caption[Comparación de metodologías ágiles móviles]{Comparación de metodologías ágiles móviles. \textit{Fuente:} Elaboración propia a partir de (Schwaber \& Sutherland, 2020); (Beck \& Andres, 2005); (Santander, 2020); (Abrahamsson, y otros, 2004).}\label{tab:metodologias-agiles}\\
\end{longtable}
}

{\small
\begin{longtable}{|L{3.6cm}|L{5.75cm}|L{5.3cm}|}
\hline
\rowcolor{encabezado} \textbf{Características} & \textbf{Android} & \textbf{iOS} \\
\hline
\endfirsthead
\hline
\rowcolor{encabezado} \textbf{Características} & \textbf{Android} & \textbf{iOS} \\
\hline
\endhead

\endfoot
\textbf{Propietario} & Google & Apple \\ \hline
\textbf{Familia S.O.} & Linux & UNIX \\ \hline
\textbf{Tipo de licencia} & Código abierto & Código cerrado y propietario \\ \hline
\textbf{Programado en} & C, C++, Java & Swift, Objective-C, C, C++ \\ \hline
\textbf{Lenguajes de programación} & Java, Kotlin & Swift \\ \hline
\textbf{Entorno de desarrollo} & Android Studio & Xcode \\ \hline
\textbf{Arquitectura del sistema} & Basada en kernel Linux + HAL + Runtime ART
& Basada en kernel XNU + frameworks de alto nivel \\ \hline
\textbf{Presencia en el mercado boliviano} & 90.74\,\% (Global, 2025) &
8.7\,\% (Global, 2025) \\ \hline
\textbf{Diseño} & Diseños flexibles, personalizado, colorido, adaptable a
muchas marcas y dispositivos. & Diseño minimalista, elegante, limpio y
extremadamente consistente. \\ \hline
\textbf{Personalización} & Permite la personalización de la interfaz mediante
launchers personalizados, permite utilizar paquetes de temas, ajustar el panel
de configuración rápida y modificar el comportamiento de aplicaciones
predeterminadas (Gironés, 2012). & Su personalización es limitada, prioriza la
simplicidad, estabilidad y coherencia visual, lo que reduce riesgos de
incompatibilidad (Apple, 2023). \\ \hline
\textbf{Fragmentación} & Samsung, Huawei, Xiaomi, Motorola, Sony. & Exclusivo
para dispositivos Apple (iPhone, iPad). \\ \hline
\textbf{Distribución de Apps} & Play Store y permite instalación manual
(APK/Sideloading). & App Store \\ \hline
\textbf{Asistente Virtual} & Google Now, Google Assistant & Siri \\ \hline
\caption[Comparación de Android vs iOS]{Comparación de Android vs iOS. \textit{Fuente:} Elaboración propia.}\label{tab:android-ios}\\
\end{longtable}
}

{\small
\begin{longtable}{|L{2.2cm}|L{5.55cm}|L{3.8cm}|L{2.7cm}|}
\hline
\rowcolor{encabezado} \textbf{Framework} & \textbf{Descripción} & \textbf{Características} &
\textbf{Ventajas} \\
\hline
\endfirsthead
\hline
\rowcolor{encabezado} \textbf{Framework} & \textbf{Descripción} & \textbf{Características} &
\textbf{Ventajas} \\
\hline
\endhead

\endfoot
React Native & Es un framework de Meta para construir aplicaciones nativas
para iOS y Android, usando JavaScript. Este framework usa el paradigma
fundamental de construcción de bloques de UI (componentes visuales con los que
interacciona el usuario) (Deloitte, 2019). &
\begin{tablista}
  \item Multiplataforma.
  \item Fast Refresh y Hot Reload.
  \item Acceso a componentes y APIs nativas.
  \item Gran comunidad y ecosistema amplio.
  \item Integración con librerías como Redux, Axios, Firebase.
\end{tablista} &
Alto rendimiento. Reutilización de código. Amplia compatibilidad. \\
\hline
Flutter & Es un framework de código abierto de Google basado en Dart con motor
de render propio (Flutter, s.f.). &
\begin{tablista}
  \item Widgets personalizables.
  \item Animaciones.
  \item Hot Reload.
\end{tablista} &
UI consistente. Excelente rendimiento gráfico. \\
\hline
Ionic & Ionic es un framework híbrido basado en tecnologías web (HTML, CSS,
JavaScript) que funciona mediante WebView. Permite construir aplicaciones con
enfoques multiplataforma usando Angular, React o Vue. Es especialmente útil
para aplicaciones orientadas a contenidos web o prototipado rápido (Ionic
Documentation, 2023). &
\begin{tablista}
  \item Multi-framework.
  \item Amplia colección de componentes UI.
  \item Soporte para Progressive Web Apps.
\end{tablista} &
Desarrollo rápido. Curva de aprendizaje baja. \\
\hline
Xamarin & Xamarin es un framework de Microsoft basado en C\# y .NET, que
compila aplicaciones nativas para Android, iOS y Windows. Permite compartir
hasta 90\,\% de la lógica de negocio entre plataformas y se integra
profundamente con Visual Studio. Es utilizado ampliamente en entornos
empresariales (Microsoft, 2022). &
\begin{tablista}
  \item Acceso directo a APIs nativas.
  \item Integración con Visual Studio.
  \item Integración con .NET MAUI.
\end{tablista} &
Sólido para apps empresariales. Reutilización de lógica. \\
\hline
Jetpack Compose & Jetpack Compose es el framework declarativo moderno de
Android, basado en Kotlin. Permite construir interfaces de usuario mediante
funciones composables, eliminando XML y ofreciendo mayor control del sistema y
rendimiento nativo. Es el estándar actual para desarrollar aplicaciones Android
modernas (Android Developers, 2024). &
\begin{tablista}
  \item UI declarativa.
  \item Integración con Android Jetpack.
  \item Animaciones avanzadas.
\end{tablista} &
Rendimiento máximo. Control total del sistema. \\ \hline
\caption[Comparación de frameworks de desarrollo Mobile]{Comparación de frameworks de desarrollo Mobile. \textit{Fuente:} Elaboración propia basada en (Deloitte, 2019); (Flutter, s.f.); (Android Developers, 2024); (Microsoft, 2022).}\label{tab:frameworks}\\
\end{longtable}
}

{\small
\begin{longtable}{|L{2.1cm}|L{5.45cm}|L{3.5cm}|L{3.2cm}|}
\hline
\rowcolor{encabezado} \textbf{Base de Datos} & \textbf{Descripción} & \textbf{Características} &
\textbf{Ventajas} \\
\hline
\endfirsthead
\hline
\rowcolor{encabezado} \textbf{Base de Datos} & \textbf{Descripción} & \textbf{Características} &
\textbf{Ventajas} \\
\hline
\endhead

\endfoot
Firebase Firestore & Es una base de datos NoSQL en la nube que permite
sincronización en tiempo real y almacenamiento distribuido (Isaac, 2025). Su
base de datos principal, Cloud Firestore, es una solución NoSQL escalable
diseñada para sincronización en tiempo real, rendimiento móvil y arquitectura
serverless. &
\begin{tablista}
  \item Escalable.
  \item Documentos y colecciones.
  \item Soporte offline.
\end{tablista} &
\begin{tablista}
  \item Fácil integración móvil.
  \item Sin servidor.
  \item Arquitectura serverless.
  \item Integración con el ecosistema de Google.
\end{tablista} \\
\hline
SQLite & Es una biblioteca en lenguaje C, es el motor de base de datos más
utilizado en el mundo. Se encuentra integrado en todos los teléfonos móviles y
en la mayoría de las computadoras, es de código abierto (SQLite Consortium,
2024). &
\begin{tablista}
  \item Escalable.
  \item Documentos y colecciones.
  \item Soporte offline.
\end{tablista} &
\begin{tablista}
  \item Fácil integración móvil.
  \item Ligera y de bajo consumo.
  \item Amplio soporte y estabilidad.
  \item Sin servidor.
  \item Tiempo real.
\end{tablista} \\
\hline
MongoDB Atlas & Es un servicio de base de datos multicloud, simplifica la
implementación y gestión de bases de datos. Permite manejar grandes volúmenes
de datos, consultas avanzadas y procesamiento en la nube con altos niveles de
seguridad y escalabilidad (MongoDB, 2024) (Huang, 2020). &
\begin{tablista}
  \item Ligero.
  \item Transaccional.
  \item Estructura relacional.
\end{tablista} &
\begin{tablista}
  \item Ideal para uso offline.
  \item Ampliamente soportado.
  \item Integración nativa con la nube.
  \item Funciones avanzadas de seguridad.
\end{tablista} \\
\hline
MySQL & MySQL es un sistema de gestión de bases de datos relacional (RDBMS) de
código abierto ampliamente utilizado en aplicaciones web y móviles.
Originalmente desarrollado por MySQL AB y actualmente mantenido por Oracle
Corporation, MySQL se ha consolidado como una solución flexible, rápida y
confiable para gestionar datos estructurados (DuBois, 2013); (Corporation,
2024). &
\begin{tablista}
  \item Lenguaje SQL para definición y manipulación de datos.
  \item Soporte para múltiples motores de almacenamiento.
  \item Ejecución de transacciones ACID.
  \item Alta compatibilidad con lenguajes de programación.
\end{tablista} &
\begin{tablista}
  \item Flexible.
  \item Escalable.
  \item Adecuada para grandes volúmenes.
\end{tablista} \\
\hline
PostgreSQL & Sistema de gestión de bases relacional ampliamente utilizado.
Destaca por su arquitectura orientada a objetos, su capacidad para manejar
grandes volúmenes de datos y su soporte para consultas complejas, extensiones y
tipos de datos personalizados. &
\begin{tablista}
  \item SQL.
  \item Robusto.
  \item Seguro.
  \item Gestión avanzada de concurrencia.
  \item Soporte para JSON y JSONB para datos semiestructurados.
\end{tablista} &
\begin{tablista}
  \item Ideal para Backend estructurado.
  \item Rendimiento estable.
  \item Perfecto para aplicaciones móviles.
\end{tablista} \\ \hline
\caption[Comparación de bases de datos para aplicaciones móviles]{Comparación de bases de datos para aplicaciones móviles. \textit{Fuente:} Elaboración basada en (Corporation, 2024); (Huang, 2020); (MongoDB, 2024).}\label{tab:bases-datos}\\
\end{longtable}
}

